%% hierarchy -- a block deletion along an edge of the reverse trie, before and after.
\begin{figure}[t]
\centering
\begin{tikzpicture}[font=\scriptsize, >={Stealth[length=1.8mm]}]
\def\h{0.34}
%% draws an upper block ending at x=#1, y=#2, with a context-dependent left part of width #3
\newcommand{\upperblock}[4]{%
  % common suffix: blocks of depths 0..b-1 (b = 9 lower blocks of width .55)
  \foreach \d in {0,...,8} {
    \pgfmathsetmacro{\xr}{#1-0.55*\d}
    \pgfmathsetmacro{\xl}{#1-0.55*(\d+1)}
    \ifnum\d>2 \ifnum\d<6
      \fill[black!22] (\xl,#2) rectangle (\xr,#2+\h);
    \fi\fi
    \draw (\xl,#2) rectangle (\xr,#2+\h);
  }
  % the part above the edge, different in each occurrence
  \pgfmathsetmacro{\xe}{#1-0.55*9}
  \draw[pattern=north east lines, pattern color=black!35] (\xe-#3,#2) rectangle (\xe,#2+\h);
  \draw[line width=1.1pt] (#1,#2-0.08) -- (#1,#2+\h+0.08);
  \draw[line width=1.1pt] (\xe-#3,#2-0.08) -- (\xe-#3,#2+\h+0.08);
  \node[left] at (\xe-#3-0.05,#2+0.5*\h) {#4};
}
%% the same upper block after the deletion: 6 lower blocks of the edge remain
\newcommand{\upperblockafter}[4]{%
  \foreach \d in {0,...,5} {
    \pgfmathsetmacro{\xr}{#1-0.55*\d}
    \pgfmathsetmacro{\xl}{#1-0.55*(\d+1)}
    \draw (\xl,#2) rectangle (\xr,#2+\h);
  }
  \pgfmathsetmacro{\xe}{#1-0.55*6}
  \pgfmathsetmacro{\xj}{#1-0.55*3}
  \draw[pattern=north east lines, pattern color=black!35] (\xe-#3,#2) rectangle (\xe,#2+\h);
  \draw[line width=1.1pt] (#1,#2-0.08) -- (#1,#2+\h+0.08);
  \draw[line width=1.1pt] (\xe-#3,#2-0.08) -- (\xe-#3,#2+\h+0.08);
  \draw[line width=1.3pt, densely dashed] (\xj,#2-0.17) -- (\xj,#2+\h+0.17);
  \node[left] at (\xe-#3-0.05,#2+0.5*\h) {#4};
}
\upperblock{9.2}{2.0}{1.3}{$w_1$:}
\upperblock{9.2}{1.2}{0.5}{$w_1$:}
\upperblock{9.2}{0.4}{2.0}{$w_2$:}
%% depth scale
\draw[->] (9.2,2.75) -- (3.9,2.75) node[left] {depth};
\foreach \d/\l in {0/0,3/i,6/j,9/b} {
  \pgfmathsetmacro{\x}{9.2-0.55*\d}
  \draw (\x,2.68) -- (\x,2.82) node[above] {$\l$};
}
\draw[densely dotted] (9.2-0.55*3,2.68) -- (9.2-0.55*3,0.3);
\draw[densely dotted] (9.2-0.55*6,2.68) -- (9.2-0.55*6,0.3);
\draw[densely dotted] (9.2-0.55*9,2.68) -- (9.2-0.55*9,0.3);
\draw[decorate, decoration={brace, mirror, amplitude=4pt}] (9.2-0.55*9,0.28) -- (9.2,0.28)
  node[midway, below=4pt] {the last $b$ lower blocks: equal in all three occurrences (the edge)};
\node[align=left, anchor=west] at (9.6,1.35) {upper anchors\\ (thick)};
\draw[->] (9.6,1.0) -- (9.25,0.75);
\node[anchor=east] at (1.6,2.75) {before:};
%% after the deletion
\begin{scope}[yshift=-3.5cm]
\upperblockafter{9.2}{2.0}{1.3}{$u_1$:}
\upperblockafter{9.2}{1.2}{0.5}{$u_1$:}
\upperblockafter{9.2}{0.4}{2.0}{$u_2$:}
\draw[decorate, decoration={brace, mirror, amplitude=4pt}] (9.2-0.55*3,0.28) -- (9.2,0.28)
  node[midway, below=4pt] {the last $i$ lower blocks};
\node[above] at (9.2-0.55*3,2.0+\h+0.17) {junctions (dashed)};
\node[anchor=east] at (1.6,2.55) {after:};
\end{scope}
\end{tikzpicture}
\caption{A block deletion at level $t$, before and after. Three upper blocks
(blocks of level $t+1$), two in $w_1$ and one in $w_2$, are cut into lower
blocks (blocks of level $t$) by the anchors of level $t$. Their parses end
with the same $b$ lower blocks, so all three traverse the same edge of the
reverse trie; they differ above it (hatched). If the cuts at the depths $j$
and $i$ carry equal labels in every traversing upper block, the lower blocks
between these two depths (gray) are deleted in all of them at once. In the
result $(u_1,u_2)$ the edited upper blocks still end with equal texts: every
junction, the place where a factor was removed, is followed by the same $i$
lower blocks.}
\label{fig:hierarchy}
\end{figure}
